\documentclass[../main.tex]{subfiles}
\begin{document}

\chapter{Gaussian Process}

\section{Definition}

\defn{GP}{}{
Consider function $f : \mathcal{X} \to  \mathbb{R}$. It is evaluated at a set of inputs
, $\boldsymbol{X} = \{ x_{n} \in  \mathcal{X} \}_{n=1}^{N}$. Let $\boldsymbol{f}_{X} = [f(\boldsymbol{x}_{1}), \ldots , f(\boldsymbol{x}_{N})]$ be the set of function values at these points.

If $\boldsymbol{f}_{X}$ is jointly Gaussian for any set of $N\ge 1$ points, then we say that $f$ is a Gaussian Process
}

We denote the GP  by 
\[
f(\boldsymbol{x}) \sim  GP(m(\boldsymbol{x}), \mathcal{K}(\boldsymbol{x}, \boldsymbol{x}'))
\]
where
\[
m(\boldsymbol{x}) = \mathbb{E}f(\boldsymbol{x}),  \quad \mathcal{K}(\boldsymbol{x}, \boldsymbol{x}') = \mathbb{E} \left[ (f(\boldsymbol{x}) - m(\boldsymbol{x}))(f(\boldsymbol{x}') - m(\boldsymbol{x}'))^{\top} \right] 
\]

For any finite set of inputs $\boldsymbol{X} = \{ \boldsymbol{x}_{1}, \ldots , \boldsymbol{x}_{N} \}$, we have
\[
p(\boldsymbol{f}_{X} | \boldsymbol{X}) = \mathcal{N}(\boldsymbol{f}_{X} | \boldsymbol{\mu}_{X}, \boldsymbol{K}_{X, X})
\]

A GP can be used to define a \emph{a prior over functions}, then update this prior with a likelihood function.


\section{GPs with Gaussian likelihoods}

\subsection{Predict using noise-free observeations}

Given $N$ data points $\mathcal{D}$, $y_{n} = f(\boldsymbol{x}_{n})$. Test data $\boldsymbol{X}_{*}$ of size $N_{*} \times  D$. We want to predict $\boldsymbol{f}_{*}(f(\boldsymbol{x}_{1}), \ldots , f(\boldsymbol{x}_{N_{*}}))$.

\[
\begin{pmatrix} 
    \boldsymbol{f}_{X} \\ 
    \boldsymbol{f}_{*}
\end{pmatrix}  \sim  \mathcal{N} \begin{pmatrix} 
    \begin{pmatrix} 
        \boldsymbol{\mu}_{X} \\ 
        \boldsymbol{\mu}_{*}
    \end{pmatrix} & \begin{pmatrix} 
        \boldsymbol{K}_{X, X} & \boldsymbol{K}_{X, *}  \\
        \boldsymbol{K}_{X,*}^{\top} & \boldsymbol{K}_{*, *}
    \end{pmatrix} 
\end{pmatrix} 
\]

\[
\begin{aligned}
p(f_{*} | \boldsymbol{X}_{*}, \mathcal{D}) &= \mathcal{N}(\boldsymbol{f}_{*} | \boldsymbol{\mu}_{* | X}, \boldsymbol{\Sigma}_{* | X}) \\
\boldsymbol{\mu}_{* | X} &= \boldsymbol{\mu}_{*} + \boldsymbol{K}_{X, *}^{\top} \boldsymbol{K}_{X, X}^{\top}(\boldsymbol{f}_{X} - \boldsymbol{\mu}_{X}) \\
\boldsymbol{\Sigma}_{* | X} &= \boldsymbol{K}_{*, *} - \boldsymbol{K}_{X, *}^{\top}\boldsymbol{K}^{-1} \boldsymbol{K}_{X, *}
\end{aligned}
\]

\subsection{Predict using noisy observations}

\[
y_{n} = f(\boldsymbol{x}_{n}) + \epsilon_{n}, \quad \operatorname{Cov}(\boldsymbol{y} | \boldsymbol{X}) = \boldsymbol{K}_{X, X} + \sigma_{y}^{2} \boldsymbol{I}_{N}.
\]


\end{document}
